% =====================================================================
% TIỂU LUẬN DẪN NHẬP CỦA BẢN CHUYỂN NGỮ
% Không thuộc nguyên tác NCTM.
% =====================================================================

\cleardoublepage
\gdef\CurrentRunningHead{Tiểu luận dẫn nhập}

% \chapter*{Tiểu luận dẫn nhập của bản chuyển ngữ}
\bookmarksetup{startatroot}
\phantomsection
\addcontentsline{toc}{chapter}{Tiểu luận dẫn nhập của bản chuyển ngữ}

\begingroup
\clubpenalty=10000
\widowpenalty=10000
\displaywidowpenalty=10000

\thispagestyle{plain}

\begin{center}
{\small\bfseries TIỂU LUẬN DẪN NHẬP CỦA BẢN CHUYỂN NGỮ\par}

\vspace{0.9\baselineskip}

{\Large\bfseries Vị trí của \emph{Focus in High School Mathematics: Reasoning and Sense Making (FHSM)}\par}

\vspace{0.15\baselineskip}

{\Large\bfseries trong giáo dục toán học đương đại\par}

\vspace{0.65\baselineskip}

{\large
Đọc từ bối cảnh Việt Nam và kỉ nguyên trí tuệ nhân tạo}
\end{center}

% \vspace{0.8\baselineskip}
\vspace{1.15\baselineskip}

\begingroup
\small

\noindent
\textit{Ghi chú về tư cách văn bản.}
Tiểu luận này do người chuyển ngữ biên soạn và không thuộc nguyên tác của Hội đồng Quốc gia Giáo viên Toán học Hoa Kì (National Council of Teachers of Mathematics, NCTM). Những nhận định về Việt Nam, về trí tuệ nhân tạo và về ý nghĩa hiện thời của cuốn sách là phần diễn giải của người viết khi đọc lại một văn bản xuất bản năm 2009. Vì vậy, chỉ những điều có căn cứ trực tiếp trong nguyên tác mới được quy cho NCTM.

\vspace{0.75\baselineskip}

\noindent
\textit{Tóm tắt.}
\emph{Focus in High School Mathematics: Reasoning and Sense Making} (NCTM,
2009) không phải là một chương trình toán trung học thay thế, một danh mục nội dung hay một sách giáo khoa. Cuốn sách đề nghị đặt suy lí và kiến tạo ý nghĩa làm nền tảng cho việc học toán. Từ điểm xuất phát đó, nó xem xét cách lựa chọn nội dung, thiết kế nhiệm vụ học tập, tổ chức dạy học, đánh giá việc học, phân phối cơ hội học tập và phối hợp giữa chương trình, dạy học với đánh giá.

Tiểu luận này trước hết xác định vị trí của FHSM trong dòng phát triển tư tưởng của NCTM và đặt cuốn sách trong đối thoại với một số định hướng của giáo dục toán học đương đại. Tiếp đó, bài viết xem FHSM có thể được đọc như một triết lí giáo dục toán học thực hành đến mức nào và đối sánh nó với Chương trình giáo dục phổ thông môn Toán hiện hành của Việt Nam. Cuối cùng, tiểu luận xem xét lại giá trị của suy lí và kiến tạo ý nghĩa trong bối cảnh trí tuệ nhân tạo tạo sinh, đặc biệt qua những hoạt động như chứng minh, diễn giải, kiểm chứng và đánh giá tính hợp lí.

Luận điểm chính của bài viết là giá trị bền vững của FHSM nằm ở việc thay đổi tiêu chuẩn dùng để nhận biết một trải nghiệm học toán có chất lượng. Không thể chỉ nhìn vào những nội dung đã được dạy, những thủ tục học sinh đã thực hiện hay việc các em có tìm được đáp án đúng hay không. Cần nhìn sâu hơn vào những hoạt động trí tuệ mà người học thực sự tham gia: các em nhận ra quan hệ và cấu trúc như thế nào, hình thành và kiểm tra phỏng đoán ra sao, dựa vào đâu để biện minh cho một kết luận, và kết nối tri thức mới với những gì đã biết như thế nào. Theo nghĩa đó, trọng tâm của FHSM không chỉ là học sinh học cái gì, mà còn là các em đang làm gì về mặt tư duy khi học toán.

Ở Việt Nam, Chương trình giáo dục phổ thông 2018 đã chính thức đặt phát triển năng lực toán học vào trung tâm. Những điều chỉnh về đánh giá, năng lực số và giáo dục trí tuệ nhân tạo sau đó tiếp tục phát triển định hướng này. Trong bối cảnh ấy, FHSM có thể đóng vai trò như một khung trung gian: nó giúp cụ thể hóa những mục tiêu năng lực tương đối khái quát thành các nhiệm vụ, tương tác và cách tổ chức dạy học trong đó học sinh thực sự phải suy lí và kiến tạo ý nghĩa trên nội dung toán học.

Sự xuất hiện của trí tuệ nhân tạo tạo sinh đặt thêm một vấn đề mà FHSM năm 2009 chưa phải đối diện. Khi công cụ có thể đảm nhận một phần đáng kể việc giải bài, tạo lập luận, viết lời giải thích hay trình bày kết quả, việc dạy học phải được thiết kế thế nào để những hoạt động nhận thức mà ta muốn phát triển vẫn thực sự được học sinh thực hiện? Đồng thời, việc đánh giá phải dựa vào những bằng chứng nào để có thể kết luận một cách đáng tin rằng năng lực đã hình thành ở chính người học? Vì vậy, FHSM nên được tiếp nhận như một nguồn tư tưởng để đối thoại, thử nghiệm và kiểm định trong những điều kiện mới, chứ không phải một mô hình ngoại nhập để sao chép.

\vspace{0.75\baselineskip}

\noindent
\textit{Từ khóa.} giáo dục toán học; suy lí (\textit{reasoning}); kiến tạo ý nghĩa (\textit{sense making}); Hội đồng Quốc gia Giáo viên Toán học Hoa Kì (NCTM); năng lực toán học; Chương trình giáo dục phổ thông 2018; giáo dục toán học Việt Nam; trí tuệ nhân tạo tạo sinh; đánh giá giáo dục; công bằng giáo dục.

\endgroup

\clearpage

\phantomsection
\addcontentsline{toc}{section}{1. Đặt vấn đề: xác định đúng loại hình và tầm vóc của FHSM}
\section*{1. Đặt vấn đề: xác định đúng loại hình và tầm vóc của FHSM}

Một cuốn sách giáo dục có thể tạo ảnh hưởng theo nhiều cách. Có cuốn thay đổi những nội dung được đưa vào chương trình; có cuốn cung cấp phương pháp dạy một chủ đề cụ thể; có cuốn tổng kết một lĩnh vực nghiên cứu. \emph{Focus in High School Mathematics: Reasoning and Sense Making}, xuất bản năm 2009, thuộc một loại khác. Chính trong phần Dẫn nhập, NCTM phân biệt ấn phẩm này với những tài liệu đương thời chủ yếu phân tích các chủ đề nên được dạy trong từng khóa toán trung học. FHSM chủ ý chuyển trọng tâm sang một câu hỏi khác: chương trình và dạy học phải được tổ chức thế nào để suy lí và kiến tạo ý nghĩa trở thành nền tảng của việc học nội dung toán học? NCTM còn mô tả ấn phẩm như một ``bộ lọc then chốt'' để xem xét cách tổ chức chương trình (NCTM 2009).

Vì vậy, sẽ không chính xác nếu xem FHSM như một chương trình thay thế cho chương trình toán trung học, càng không nên đọc nó như một tuyển tập kĩ thuật sư phạm. Câu hỏi mà nó đặt ra nằm ở một tầng khác: bất kể một hệ thống chọn dạy những chủ đề nào, học sinh đang được yêu cầu \emph{làm gì về mặt trí tuệ} khi học những chủ đề ấy? Các em có nhận ra quan hệ và cấu trúc, hình thành phỏng đoán, tìm căn cứ, kiểm tra tính hợp lí, giải thích và biện minh, đánh giá lập luận, diễn giải kết quả trong bối cảnh, kết nối tri thức mới với tri thức đã có hay không? Hay hoạt động học tập chủ yếu dừng ở việc tiếp nhận thủ tục, nhận dạng dạng vấn đề và tái tạo thao tác trong những tình huống quen thuộc?

Cách đặt vấn đề này cho thấy tầm vóc thật sự của cuốn sách. Tầm vóc ấy không nằm ở chỗ các ví dụ sử dụng kiến thức toán học cao đến đâu; phần lớn toán học trong FHSM vẫn thuộc phạm vi phổ thông. Cuốn sách cũng không xây dựng một hệ triết học hoàn chỉnh. Điều quan trọng là FHSM muốn thay đổi \emph{tiêu chuẩn để nhìn vào việc học}. Nếu suy lí và kiến tạo ý nghĩa là nền tảng chứ không phải phần thêm vào sau cùng, ta không thể chỉ thay vài bài tập. Cách thiết kế nhiệm vụ học tập, tổ chức đối thoại trong lớp, sử dụng công nghệ, hiểu về chứng minh, đánh giá việc học, phân phối cơ hội học tập và phối hợp chương trình--dạy học--đánh giá đều phải được xem xét lại dưới cùng tiêu chuẩn ấy.

Cấu trúc ba phần của FHSM phản ánh đúng tham vọng ấy. Phần thứ nhất xác định suy lí và kiến tạo ý nghĩa, mô tả các thói quen suy lí và lộ trình phát triển của chúng. Phần thứ hai đưa hai quá trình này vào năm vùng nội dung lớn của toán trung học: số và đo lường, kí hiệu đại số, hàm số, hình học, thống kê và xác suất. Phần thứ ba rời khỏi phạm vi một bài học đơn lẻ để xem xét công bằng, tính nhất quán giữa chương trình--dạy học--đánh giá, và trách nhiệm của học sinh, gia đình, giáo viên, nhà quản lí, nhà hoạch định chính sách, giáo dục đại học và những người phát triển chương  rình. Việc cuốn sách bắt đầu từ hoạt động nhận thức của học sinh nhưng kết thúc ở cấu trúc của cả hệ thống giáo dục cho thấy phạm vi của nó rộng hơn ``phương pháp dạy học''. FHSM quan tâm tới những điều kiện cần có để một hệ thống giáo dục tạo cơ hội cho học sinh hình thành một kiểu hiểu biết toán học dựa trên suy lí và kiến tạo ý nghĩa.

\phantomsection
\addcontentsline{toc}{section}{2. FHSM trong dòng phát triển tư tưởng của NCTM}
\section*{2. FHSM trong dòng phát triển tư tưởng của NCTM}

FHSM không xuất hiện như một sáng kiến biệt lập. Chính phần Dẫn nhập đặt cuốn sách vào một tiến trình dài trong các văn kiện NCTM. \emph{An Agenda for Action} (1980) đưa giải quyết vấn đề vào vị trí nổi bật. \emph{Curriculum and Evaluation Standards for School Mathematics} (1989) nhấn mạnh một lõi toán học chung cùng các quá trình giải quyết vấn đề, suy lí, kết nối và giao tiếp. \emph{Principles and Standards for School Mathematics} (2000) tổ chức chương trình bằng năm Chuẩn nội dung và năm Chuẩn quá trình. \emph{Curriculum Focal Points} (2006) tiếp tục tìm cách làm chương trình tập trung và nhất quán hơn ở các lớp trước trung học, nhưng vẫn yêu cầu các quá trình toán học phải hiện diện trong chính việc học nội dung. Có thể xem FHSM như một bước nhấn mạnh tiếp theo của dòng phát triển ấy: cuốn sách không thêm một danh mục nội dung mới mà đưa suy lí và kiến tạo ý nghĩa vào vị trí nền tảng cho những nội dung và quá trình đã được xác lập trước đó (NCTM 1980, 1989, 2000, 2006, 2009).

Sự tập trung vào suy lí cũng không chỉ là câu chuyện nội bộ của NCTM. Khung Toán học của Chương trình Đánh giá Học sinh Quốc tế (Programme for International Student Assessment, PISA) năm 2022, do Tổ chức Hợp tác và Phát triển Kinh tế (Organisation for Economic Co-operation and Development, OECD) xây dựng, đặt suy lí toán học (\textit{mathematical reasoning}) vào vị trí trung tâm của năng lực toán học. Bên cạnh suy lí, khung này còn xét ba quá trình: hình thành một vấn đề dưới dạng toán học, sử dụng toán học, và diễn giải--đánh giá kết quả trong bối cảnh. OECD lí giải sự điều chỉnh đó bằng một thế giới thay đổi nhanh dưới tác động của công nghệ, nơi công dân phải tự đưa ra phán đoán và sử dụng toán học trong những tình huống không quen thuộc (OECD 2023). Sự gần gũi này không chứng minh FHSM đã ảnh hưởng trực tiếp tới PISA. Nó chỉ cho thấy vấn đề mà FHSM đặt ở trung tâm vẫn là một mối quan tâm quan trọng của giáo dục toán học đương đại.

Điểm này quan trọng vì nó ngăn một cách đọc giản lược: ``NCTM chuyển từ dạy kiến thức sang dạy tư duy''. FHSM không phủ định nội dung toán học và cũng không đối lập năng lực với tri thức. Chương 1 nói rõ trọng tâm vào suy lí và kiến tạo ý nghĩa phải được phát triển \emph{trong bối cảnh của nội dung quan trọng}; học sinh vẫn phải thực hiện chính xác các thủ tục, nhưng đồng thời phải hiểu vì sao chúng có hiệu lực, biết khi nào có thể dùng chúng và biết diễn giải kết quả. Quan điểm này gần với cấu trúc ``năng lực toán học'' trong \emph{Adding It Up}, nơi sự hiểu biết khái niệm, sự thành thạo thủ tục, năng lực chiến lược, suy lí thích ứng và khuynh hướng tích cực được xem như những mạch đan xen của một năng lực thống nhất (Kilpatrick, Swafford và Findell 2001).

Bởi thế, FHSM không buộc ta chọn một phía giữa ``kiến thức'' với ``năng lực'', hay giữa ``kĩ năng'' với ``hiểu''. Vấn đề mà cuốn sách đặt ra là kiến thức được hình thành trong đầu người học theo cấu trúc nào. Một thủ tục có thể chỉ được nhớ như chuỗi thao tác phải làm lần lượt. Nhưng cũng thủ tục ấy có thể được gắn với biểu diễn, khái niệm, điều kiện áp dụng và lí do khiến nó đúng. Vì vậy, hai học sinh có thể cho cùng một đáp số nhưng hiểu toán học theo hai cách rất khác nhau.

Ví dụ về công thức khoảng cách ở Chương 1 thể hiện sự khác biệt này với độ tương phản đặc biệt rõ. Trong một kịch bản, học sinh cố khôi phục công thức từ những mảnh kí hiệu đã nhớ, thậm chí lẫn với công thức trung điểm. Trong kịch bản khác, bài toán được tái dựng từ tình huống hình học: sự thay đổi theo hai phương vuông góc dẫn tới tam giác vuông, định lí Pythagoras cung cấp quan hệ, và công thức tổng quát xuất hiện như kết quả của một chuỗi kết nối. Điều đáng chú ý không phải giáo viên đã ``giải thích hay hơn'', mà là học sinh được tham gia vào một loại hoạt động khác: các em có thể tái tạo công thức vì công thức đã được đặt trong một mạng lưới lí do và quan hệ. FHSM dùng ví dụ ấy để cho thấy suy lí và kiến tạo ý nghĩa vừa là kết quả của dạy học, vừa là phương tiện để học sinh đi đến hiểu biết về toán học (NCTM 2009).

Từ đây có thể hiểu vì sao Chương 2 nhấn mạnh rằng suy lí không phải một nhóm chủ đề mới để thêm vào chương trình. Nó là một \emph{định hướng xuyên suốt đối với việc học toán}. Bởi vậy, câu hỏi quan trọng không phải ``đã dành bao nhiêu tiết cho suy lí''. Cần hỏi trong mỗi miền nội dung, học sinh có thường xuyên phải tìm cấu trúc, đưa ra và kiểm tra phỏng đoán, lựa chọn biểu diễn, xây dựng căn cứ, xem lại lời giải và khái quát hóa hay không. Nói cách khác, suy lí không phải một ``nội dung thứ sáu'' đứng cạnh đại số, hình học hay thống kê; nó là một cách hoạt động toán học xuyên qua những nội dung ấy.

\phantomsection
\addcontentsline{toc}{section}{3. Một triết lí giáo dục toán học ở mức nào?}
\section*{3. Một triết lí giáo dục toán học ở mức nào?}

Câu hỏi liệu FHSM có chứa một triết lí giáo dục hay không cần được trả lời với sự dè dặt về thuật ngữ. Nếu ``triết lí giáo dục'' được hiểu là một hệ thống hoàn chỉnh về bản chất của tri thức, con người, giá trị, mục đích tối hậu của giáo dục và trật tự xã hội, FHSM rõ ràng không phải một chuyên luận như thế. NCTM không xây dựng một bản thể luận về toán học, không khảo cứu có hệ thống
các trường phái nhận thức luận, cũng không trình bày một triết học chính trị của giáo dục.

Nhưng có thể dùng ``triết lí giáo dục'' theo một nghĩa hẹp và thực hành hơn. Khi ấy, ta hỏi một văn bản có đưa ra những quan niệm tương đối nhất quán về việc biết toán là gì, người học giữ vai trò nào trong việc hình thành sự hiểu biết, hoạt động toán học nào được coi là có giá trị, và hệ thống giáo dục phải tạo những điều kiện gì để các giá trị ấy trở thành hiện thực hay không. Theo nghĩa này, FHSM cho phép một cách đọc có chiều sâu triết lí. Cần nhấn mạnh rằng đây là \emph{phép phân tích của tiểu luận}; NCTM không tự gọi FHSM là một hệ triết học.

Trước hết là một quan niệm về \emph{biết toán}. Đối với FHSM, biết toán không thể chỉ là nhớ một tập hợp kết quả và thủ tục. Suy lí được hiểu rộng là quá trình đi tới kết luận dựa trên chứng cứ hoặc những giả định đã nêu. Trong toán học, quá trình ấy có thể bắt đầu từ quan sát quy nạp, tìm tòi, phỏng đoán, giải thích và biện minh phi hình thức; nó cũng có thể phát triển tới suy diễn hình thức và chứng minh. Kiến tạo ý nghĩa là phát triển sự hiểu biết về một tình huống, bối cảnh hay khái niệm bằng cách nối cái mới với tri thức đã có. Vì vậy, tri thức toán học có chất lượng không chỉ trả lời ``cái gì'', mà còn giúp người học trả lời ``vì sao'', ``liên hệ với cái gì'', ``dựa vào đâu'', ``đúng trong điều kiện nào'' và ``có nghĩa gì trong bối cảnh''.

Tiếp đó là một quan niệm về \emph{vai trò của người học}. FHSM nhiều lần nhấn mạnh rằng học sinh phải tự mình tham gia vào suy lí, chứ không chỉ quan sát giáo viên hay sách giáo khoa suy lí thay mình. Điều này dẫn tới một hệ quả quan trọng: học sinh không thể chỉ tiếp nhận những tri thức đã được hoàn tất sẵn. Các em phải nêu căn cứ, đánh giá suy lí của mình và của người khác, lựa chọn
biểu diễn, kiểm tra phỏng đoán và chịu trách nhiệm về những kết luận mình chấp nhận.

Trong tiểu luận này, phẩm chất đó được gọi là \emph{tư cách chủ thể nhận thức}. Cụm từ này chỉ một điều khá cụ thể: người học không đứng ngoài quá trình xác lập điều gì là đúng, mà tự tham gia vào việc tìm căn cứ, kiểm tra và quyết định một khẳng định toán học có đáng chấp nhận hay không. Đây là công cụ diễn giải của người viết, không phải thuật ngữ mà FHSM tự sử dụng.

FHSM cũng cho phép xem lại quan hệ giữa \emph{phi hình thức và hình thức}. Cuốn sách không chia việc học thành hai phía tách biệt, trong đó kiến tạo ý nghĩa chỉ thuộc về trực giác còn suy lí chỉ bắt đầu khi xuất hiện chứng minh hình thức. Chương 1 đặt hai quá trình này trong quan hệ đan xen. Chứng minh là một hình thức suy lí có mức độ chặt chẽ cao, nhưng nó được xây dựng trên nền của việc hiểu vấn đề và các quan hệ toán học. Ngược lại, một chứng minh tốt có thể làm lộ cấu trúc, giải thích vì sao kết quả đúng và giúp người học kiểm tra tính hợp lệ của suy lí.

Từ đó có thể rút ra một nguyên tắc giáo dục: khi đi tới những hình thức toán học chặt chẽ hơn, không nên đánh mất ý nghĩa. Nhưng trực giác và giải thích phi hình thức cũng không phải điểm dừng cuối cùng nếu vấn đề đòi hỏi một chứng minh chặt chẽ hơn.

Cuối cùng là một quan niệm về \emph{công bằng và thể chế}. FHSM dành hẳn Chương 9 cho công bằng, nhưng công bằng ở đây không chỉ có nghĩa mọi học sinh được tiếp cận cùng một nội dung. Còn phải hỏi các em có được tiếp cận cùng một chất lượng hoạt động toán học hay không. Nếu một nhóm thường xuyên được làm những nhiệm vụ yêu cầu phỏng đoán, giải thích, chứng minh và khái quát, trong khi nhóm khác chủ yếu luyện thủ tục vì bị cho là ``chưa đủ khả năng'', thì cơ hội học tập đã được phân phối không đều.

Chương 10 mở rộng câu hỏi ấy sang sự nhất quán giữa chương trình, dạy học và đánh giá; Chương 11 tiếp tục đặt trách nhiệm lên nhiều bên liên quan. Vì vậy, FHSM không chỉ đòi hỏi từng giáo viên thay đổi cách dạy. Cuốn sách còn quan tâm tới những điều kiện của cả hệ thống có thể duy trì cơ hội tham gia vào hoạt động toán học có chất lượng cho mọi học sinh.

Trong phạm vi nói trên, có thể gọi FHSM là một \emph{triết lí giáo dục toán học thực hành}. Cách gọi này muốn nói rằng cuốn sách đưa ra một quan niệm có tính chuẩn tắc về thế nào là học toán có chất lượng, rồi theo đuổi quan niệm ấy từ hoạt động của từng học sinh tới cách tổ chức của cả hệ thống. Đây chỉ là một nhãn phân tích hữu ích; nó không có nghĩa FHSM là một triết thuyết hoàn chỉnh. Cuốn sách được viết cho toán trung học và cho những vấn đề của giáo dục Hoa Kì ở cuối thập niên 2000; nó không phải một lí thuyết đầy đủ về giáo dục toán học từ mầm non tới nghiên cứu toán học, càng không phải triết lí duy nhất mà một hệ thống giáo dục phải theo.

\phantomsection
\addcontentsline{toc}{section}{4. Giáo dục toán học Việt Nam: từ mục tiêu năng lực tới câu hỏi về hoạt động học}
\section*{4. Giáo dục toán học Việt Nam: từ mục tiêu năng lực tới câu hỏi về hoạt động học}

Đặt FHSM vào Việt Nam không nhất thiết phải bắt đầu bằng một lịch sử dài của toán học hay giáo dục toán học Việt Nam. Cũng không nên mở đầu bằng những nhận định rộng như lớp học Việt Nam chỉ có học thuộc, luyện thi hay thao tác thủ tục. Muốn nói như vậy cần có bằng chứng thực chứng đủ rộng và đủ đáng tin.

Với câu hỏi của tiểu luận này, một điểm xuất phát chắc chắn hơn là những lựa chọn đã được Việt Nam chính thức ghi vào chương trình và hệ thống đánh giá. Từ đó mới đặt câu hỏi: FHSM có thể giúp những mục tiêu ấy trở thành hoạt động học cụ thể trong lớp Toán hay không?

Chương trình giáo dục phổ thông 2018 tạo ra một chuyển hướng rõ ở cấp mục tiêu chính thức. Giáo dục toán học được đặt trong định hướng phát triển phẩm chất, năng lực chung và năng lực toán học. Năm thành phần cốt lõi của năng lực toán học gồm: tư duy và lập luận toán học; mô hình hóa toán học; giải quyết vấn đề toán học; giao tiếp toán học; sử dụng công cụ, phương tiện học toán.

Tuy vậy, chương trình không vì thế hạ thấp kiến thức và kĩ năng. Nó vẫn quy định những kiến thức, kĩ năng toán học then chốt và tổ chức nội dung quanh các mạch Số, Đại số và Một số yếu tố giải tích; Hình học và Đo lường; Thống kê và Xác suất. Chương trình cũng nhấn mạnh kết nối giữa các ý tưởng toán học, giữa toán học với thực tiễn và với những lĩnh vực khác (Bộ Giáo dục và Đào tạo 2018). Vì vậy, ``năng lực'' và ``nội dung'' không đứng ở hai phía đối lập. Nội dung chính là môi trường trong đó năng lực phải được hình thành.

Vì thế, khi đọc FHSM ở Việt Nam, câu hỏi không còn là giáo dục toán học có nên ``chuyển sang năng lực'' hay không. Ở cấp chương trình quốc gia, lựa chọn ấy đã được xác lập. Câu hỏi khó hơn nằm ở bước tiếp theo. Trong một bài học cụ thể, ``tư duy'', ``lập luận'', ``mô hình hóa'', ``giải quyết vấn đề'' và ``giao tiếp'' phải hiện ra thành những hành động toán học nào? Những hành động ấy được nuôi dưỡng qua nội dung ra sao? Và giáo viên dựa vào bằng chứng nào để biết năng lực đang thực sự hình thành ở người học?

Những thay đổi về đánh giá cho thấy định hướng phát triển năng lực không chỉ dừng ở văn bản chương trình. Từ kì thi tốt nghiệp trung học phổ thông năm 2025, Bộ Giáo dục và Đào tạo đã áp dụng cấu trúc đề thi mới theo hướng phù hợp với đánh giá năng lực của Chương trình giáo dục phổ thông 2018, đồng thời bổ sung những dạng câu hỏi mới bên cạnh trắc nghiệm nhiều lựa chọn truyền thống (Cục Quản lí chất lượng 2023). Đến năm 2026, quy định về đánh giá diện rộng cấp quốc gia tiếp tục nhấn mạnh việc đánh giá phẩm chất và năng lực của học sinh (Bộ Giáo dục và Đào tạo 2026a). Những thay đổi này cho thấy hướng điều chỉnh của hệ thống đánh giá chính thức. Tuy nhiên, từ các văn bản chính sách và cấu trúc đề thi không thể suy ra rằng những năng lực mà chương trình đặt ra đã được đánh giá đầy đủ, càng không thể suy ra rằng chúng đã được hình thành đầy đủ trong mọi lớp học.

Bối cảnh công nghệ cũng đang được bổ sung trực tiếp vào hệ thống năng lực. Khung năng lực số cho người học ban hành năm 2025 dành một trong sáu miền năng lực cho ứng dụng trí tuệ nhân tạo, nhấn mạnh việc hiểu, sử dụng và đánh giá hệ thống AI một cách có đạo đức và trách nhiệm (Bộ Giáo dục và Đào tạo 2025). Đến tháng 8 năm 2026, Bộ Giáo dục và Đào tạo ban hành Khung nội dung giáo dục trí tuệ nhân tạo cho học sinh phổ thông, triển khai đại trà từ năm học 2026--2027. Khung mới yêu cầu rèn luyện tư duy phản biện, khả năng kiểm chứng thông tin và sử dụng AI an toàn, có đạo đức, có trách nhiệm. Đồng thời, việc triển khai không được làm thay đổi mục tiêu hay gia tăng yêu cầu cần đạt của Chương trình giáo dục phổ thông; nội dung AI có thể được lồng ghép trong các môn học (Bộ Giáo dục và Đào tạo 2026b, 2026c). Điều này cho thấy giáo dục AI đang được bổ sung vào Chương trình 2018 mà không thay thế định hướng phát triển năng lực đã có. Nói cách khác, AI làm xuất hiện những yêu cầu mới, nhưng các yêu cầu ấy được đặt vào bên trong cấu trúc chương trình hiện hành chứ không xây một chương trình song song.

Trong cấu trúc ấy, FHSM không cần và cũng không nên cạnh tranh với chương trình quốc gia. Quan hệ thích hợp hơn là quan hệ bổ trợ. Chương trình Việt Nam xác định học sinh cần phát triển những năng lực nào và phải đạt những yêu cầu gì. FHSM có thể giúp cụ thể hóa một câu hỏi khác: trong quá trình học những nội dung ấy, học sinh cần được đặt vào những nhiệm vụ và tương tác nào để chính các em phải suy lí và kiến tạo ý nghĩa? Chẳng hạn, các em đang tìm quan hệ nào, dựa trên căn cứ nào, hình thành và kiểm tra phỏng đoán ra sao, lựa chọn biểu diễn thế nào, diễn giải kết quả trong bối cảnh nào, và xem lại tính hợp lí của lời giải đến mức nào.

Đây là tầng của \emph{hoạt động học}. Chính ở tầng này, một văn bản nước ngoài có thể hữu ích cho chương trình Việt Nam mà không cần trở thành mô hình để sao chép. FHSM không thay chương trình xác định đích đến; nó giúp đặt câu hỏi cụ thể hơn về những hoạt động toán học mà người học phải thực sự tham gia trên đường đi tới những đích ấy.

\phantomsection
\addcontentsline{toc}{section}{5. Kỉ nguyên trí tuệ nhân tạo: hoạt động học, bằng chứng và năng lực}
\section*{5. Kỉ nguyên trí tuệ nhân tạo: hoạt động học, bằng chứng và năng lực}

FHSM được xuất bản năm 2009. Mọi diễn giải cuốn sách trong bối cảnh trí tuệ nhân tạo tạo sinh phải bắt đầu bằng sự thật lịch sử ấy. NCTM không thể viết FHSM như một phản ứng đối với những hệ thống tạo sinh phổ biến từ đầu thập niên 2020; vì vậy, không nên gán cho văn bản những tiên đoán mà nó chưa từng đưa ra. Điều có thể làm một cách nghiêm túc là hỏi: những thay đổi công nghệ hiện nay làm biến đổi điều kiện của việc dạy, học và đánh giá toán học ra sao, và trong những điều kiện mới ấy, các luận điểm của FHSM còn giữ giá trị nào và cần được bổ sung ở đâu?

Điểm mới của trí tuệ nhân tạo tạo sinh không phải ở chỗ máy móc lần đầu tiên có thể làm toán. Máy tính bỏ túi, hệ đại số máy tính, phần mềm hình học động và công cụ thống kê từ lâu đã thực hiện thay con người một phần thao tác. Điều mới hơn là AI có thể tham gia sâu vào những hoạt động vốn thường được giao cho người học: đề xuất chiến lược, giải bài, tạo lập luận, viết lời giải thích,
diễn giải kết quả, viết mã hoặc trình bày một văn bản mang hình thức chứng minh. Vì thế, AI có thể làm thay không chỉ một số thao tác tính toán mà cả một phần công việc nhận thức mà một nhiệm vụ học tập vốn được thiết kế để học sinh tự thực hiện.

Đây trước hết là một vấn đề của \emph{thiết kế việc học}. Nếu suy lí và kiến tạo ý nghĩa là những quá trình nền tảng của việc học toán, như FHSM chủ trương, thì một nhiệm vụ chỉ có giá trị đối với mục tiêu ấy khi học sinh thực sự phải tham gia vào các quá trình đó. Việc một công cụ giúp các em làm được nhiều hơn chưa tự nó cho biết các em đã học được nhiều hơn. Cần hỏi phần nào của việc phân tích vấn đề, tìm quan hệ, lựa chọn biểu diễn, hình thành và kiểm tra phỏng đoán, biện minh hay diễn giải vẫn do người học thực hiện, và phần nào đã được giao cho máy.

Một thí nghiệm thực địa của Bastani và cộng sự (2025), thực hiện với gần một nghìn học sinh trung học trong các giờ học toán, cho thấy vì sao cần phân biệt giữa việc làm bài tốt hơn khi có AI hỗ trợ với việc người học thực sự phát triển được năng lực của mình. Trong nghiên cứu này, học sinh được sử dụng những công cụ dựa trên GPT-4 để hỗ trợ giải bài. Khi công cụ còn được phép sử dụng, các em có thể đạt kết quả tốt hơn so với khi không có hỗ trợ.

Tuy nhiên, điều quan trọng xuất hiện khi công cụ được rút đi và học sinh phải tự làm bài. Với GPT Base, một phiên bản ít được thiết kế để hạn chế việc cung cấp trực tiếp lời giải, những học sinh trước đó đã làm bài tốt hơn nhờ AI lại có kết quả thấp hơn nhóm đối chứng khi phải làm việc độc lập. Nói cách khác, thành tích cao trong lúc có công cụ không nhất thiết đồng nghĩa với việc những hiểu biết và năng lực cần thiết đã được hình thành ở chính người học.

Kết quả lại khác với GPT Tutor. Công cụ này được thiết kế theo hướng không làm thay quá nhiều phần việc mà học sinh cần tự thực hiện, chẳng hạn bằng cách hướng dẫn từng bước thay vì đơn giản cung cấp câu trả lời. Khi sử dụng thiết kế như vậy, tác động bất lợi đối với kết quả học tập sau khi rút công cụ giảm đáng kể.

Nghiên cứu này vì thế không cho phép kết luận rằng AI nói chung có lợi hay có hại cho việc học toán. Điều nó cho thấy rõ hơn là hai việc phải được phân biệt: một công cụ có thể giúp học sinh thực hiện tốt hơn nhiệm vụ trước mắt, nhưng điều đó chưa bảo đảm rằng học sinh đã học được những gì cần thiết để sau đó tự mình thực hiện nhiệm vụ. Việc AI hỗ trợ hay làm suy yếu sự học phụ thuộc một phần quan trọng vào cách công cụ được thiết kế và vào việc những hoạt động nhận thức nào vẫn được giữ lại cho người học tự thực hiện.

AI đồng thời đặt ra một vấn đề khác ở \emph{đánh giá}. Một bài làm, một lời giải thích hay một cách trình bày quan sát được không phải chính năng lực của người học; chúng là những biểu hiện từ đó giáo viên hoặc một hệ thống đánh giá suy ra điều người học biết và có thể làm. \emph{Knowing What Students Know} của National Research Council mô tả đánh giá như một quá trình suy luận từ bằng chứng, trong đó ba thành phần phải gắn với nhau: một quan niệm về tri thức và năng lực cần đánh giá; những nhiệm vụ hay tình huống tạo ra điều có thể quan sát; và cách diễn giải các quan sát ấy để đưa ra nhận định về người học (National Research Council 2001). Theo cách nhìn này, một câu trả lời đúng có thể là bằng chứng hữu ích, nhưng không bao giờ tự nó đồng nhất với năng lực mà ta muốn kết luận.

Trí tuệ nhân tạo làm cho bước suy luận từ bằng chứng ấy trở nên khó hơn trong một số tình huống. Nếu AI đã góp phần đáng kể vào lời giải, lời giải thích hay lập luận mà ta quan sát, thì cần biết điều gì trong đó phản ánh hoạt động của người học và điều gì phản ánh năng lực của công cụ. Vấn đề không phải là mọi bài làm có AI đều mất giá trị đánh giá. Vấn đề là điều kiện sử dụng công cụ phải được tính vào khi diễn giải bằng chứng. Một nhiệm vụ dùng để đánh giá năng lực độc lập của học sinh đòi hỏi điều kiện khác với một nhiệm vụ đánh giá khả năng sử dụng AI để giải quyết vấn đề, kiểm chứng kết quả hoặc so sánh nhiều chiến lược.

Lập trường NCTM sau khi trí tuệ nhân tạo tạo sinh trở nên phổ biến đi theo một hướng tương thích với cách đặt vấn đề này. Trong tuyên bố \emph{Artificial Intelligence and Mathematics Teaching} năm 2024, NCTM khẳng định công cụ AI không thay thế nhu cầu dạy toán và giải quyết vấn đề; khả năng tạo ra câu trả lời sai hoặc không hợp lí khiến kiến thức nền và trực giác toán học vẫn cần thiết để thẩm định đầu ra. NCTM cũng xem sự tồn tại của các công cụ giải nhanh như một áp lực để tránh những đánh giá nông, chỉ đo thao tác, và chuyển một phần trọng tâm từ ``giải'' sang sự kết hợp giữa ``giải'' với ``kiểm chứng''; việc kiểm chứng nhiều lời giải có thể đòi hỏi hiểu biết sâu hơn việc chỉ tạo ra một lời giải (NCTM 2024).

UNESCO đặt vấn đề ở phạm vi rộng hơn. Hướng dẫn về trí tuệ nhân tạo tạo sinh trong giáo dục và nghiên cứu yêu cầu việc dùng AI phải bảo vệ \emph{tính chủ thể của con người} (\emph{human agency}). Nói đơn giản, người học vẫn phải là người hiểu mình đang làm gì, đưa ra những quyết định quan trọng và chịu trách nhiệm về việc sử dụng kết quả, thay vì chỉ làm theo đầu ra của máy. UNESCO đồng thời nhấn mạnh sự tham gia tích cực, tư duy bậc cao và trách nhiệm của con người đối với độ chính xác của nội dung do AI tạo ra (Miao và Holmes 2023). Khung năng lực AI dành cho học sinh tiếp tục yêu cầu người học biết đánh giá phê phán các giải pháp AI và hiểu ranh giới giữa vai trò của con người với vai trò của máy (Miao, Shiohira và Lao 2024). Những thảo luận sau đó của UNESCO mở rộng vấn đề sang thiết kế đánh giá, bất bình đẳng tiếp cận, quyền của người học và các hình thức đồng kiến tạo người--máy, nhưng vẫn đặt quyền con người, công bằng và sự bao trùm ở trung tâm (UNESCO 2025a, 2025b).

Như vậy, phần giao nhau giữa FHSM với vấn đề AI không nằm ở việc xem suy lí và kiến tạo ý nghĩa như những thuộc tính phải ``tìm thấy'' trong một sản phẩm. Hai khái niệm này vẫn giữ đúng vị trí của chúng trong FHSM: đó là những quá trình người học cần thực sự tham gia khi học toán. AI đặt thêm hai câu hỏi bao quanh các quá trình ấy. Về dạy học: việc dùng công cụ có còn để lại cho học sinh những hoạt động nhận thức mà bài học muốn phát triển hay không? Về đánh giá: những gì ta quan sát được, trong điều kiện có hoặc không có công cụ, cho phép ta suy ra điều gì về hiểu biết và năng lực của chính người học?

Từ đó, không thể giải quyết vấn đề chỉ bằng hai lựa chọn ``cấm AI'' hoặc ``cho phép AI''. Trước hết phải xác định mục tiêu của hoạt động. Có lúc cần học sinh tự thực hiện một quá trình để hình thành hoặc bộc lộ năng lực độc lập, nên công cụ phải được giới hạn. Có lúc khả năng sử dụng công cụ một cách có hiểu biết lại là một phần của năng lực cần phát triển và đánh giá. Trong những hoạt động khác, học sinh có thể được dùng AI nhưng phải thẩm định đầu ra, so sánh chiến lược, truy tìm lỗi, nêu điều kiện áp dụng, xây dựng phản ví dụ hoặc giải thích vì sao một kết quả đáng được chấp nhận.

AI buộc giáo dục phải xác định rõ hơn hoạt động nào người học cần tự thực hiện, công cụ được phép hỗ trợ ở đâu, và bằng chứng nào đủ sức nâng đỡ một kết luận về năng lực. Theo nghĩa đó, AI không làm FHSM lỗi thời; nó làm nổi bật hơn yêu cầu cốt lõi của FHSM rằng người học phải thực sự tham gia vào hoạt động toán học có ý nghĩa. Đồng thời, những vấn đề về tính chủ thể, trách nhiệm, kiểm chứng, dữ liệu, quyền riêng tư, thiên lệch và quyền của người học cho thấy một văn bản năm 2009 không thể tự mình cung cấp một khung hoàn chỉnh cho giáo dục toán học trong thời đại AI.

\phantomsection
\addcontentsline{toc}{section}{6. Tầm nhìn xã hội của FHSM: công dân, công việc, khoa học và công bằng}
\section*{6. Tầm nhìn xã hội của FHSM: công dân, công việc, khoa học và công bằng}

Một cách đọc chỉ tập trung vào phương pháp lớp học sẽ bỏ qua chiều xã hội của FHSM. Ngay đầu Chương 1, NCTM đặt toán trung học trong ba miền chuẩn bị rộng: toán học cho đời sống, cho nơi làm việc, và cho cộng đồng khoa học--kĩ thuật. Từ đó, suy lí và kiến tạo ý nghĩa không chỉ được biện minh bằng việc giúp học sinh làm bài tốt hơn. NCTM liên hệ chúng với khả năng đưa ra quyết định có căn cứ trong đời sống cá nhân và công cộng, với khả năng thích ứng trong công việc, và với việc tiếp tục học ở những lĩnh vực khoa học và kĩ thuật (NCTM 2009).

Kỉ nguyên dữ liệu và trí tuệ nhân tạo làm chiều xã hội này rõ hơn, đồng thời phức tạp hơn. Công dân hiện đại không phải lúc nào cũng tự thực hiện mọi phép tính hay tự tạo ra mọi phân tích. Thường xuyên hơn, họ phải đọc kết quả do người khác hoặc một hệ thống tạo ra rồi tự quyết định có nên tin kết quả ấy hay không. Muốn làm được vậy, họ phải hiểu đại lượng đang đo điều gì, mô hình dựa trên giả định nào, tương quan có đồng nghĩa với quan hệ nhân quả hay không, mức bất định lớn đến đâu và kết luận chỉ có hiệu lực trong phạm vi nào. Họ cũng cần nhận ra khi một con số được trình bày với độ chính xác không hợp lí.

Trong môi trường AI, còn phải thêm một năng lực nữa: nhận ra một câu trả lời có thể rất trôi chảy và thuyết phục nhưng vẫn sai hoặc thiếu căn cứ, rồi biết cách kiểm chứng nó. Những hoạt động này gặp trực tiếp các thói quen suy lí mà FHSM nhấn mạnh, đặc biệt trong thống kê và xác suất.

Chiều xã hội ấy còn thể hiện ở quan niệm về công bằng. Nếu suy lí là năng lực cần cho đời sống công dân, nghề nghiệp và việc học tiếp tục, thì nó không thể là đặc quyền của học sinh ở các lớp nâng cao. Chương 9 của FHSM bàn tới phân luồng, nhân khẩu học, kì vọng, niềm tin và thiên kiến vì những yếu tố này có thể ảnh hưởng tới loại cơ hội toán học mà học sinh được nhận.

Theo cách nhìn ấy, ``công bằng trong giáo dục toán học'' không chỉ là được đến trường hay được học cùng một nội dung. Còn phải hỏi học sinh có được làm những nhiệm vụ đòi hỏi tư duy, có được tham gia những cuộc đối thoại trong đó lí lẽ được coi trọng, và có được hỗ trợ để tiến từ suy lí phi hình thức tới những hình thức chặt chẽ hơn hay không.

Khi AI tham gia vào giáo dục, câu hỏi công bằng cần được đặt thêm ở hai tầng. Thứ nhất là khả năng tiếp cận: UNESCO cảnh báo sự chênh lệch về hạ tầng, thuê bao, ngôn ngữ và dữ liệu có thể biến khoảng cách số thành một khoảng cách AI mới (UNESCO 2025a, 2025b). Thứ hai là chất lượng của hoạt động: ngay cả khi hai học sinh cùng có công cụ, một em có thể được hướng dẫn dùng AI để đặt câu hỏi, kiểm chứng và mở rộng suy nghĩ, trong khi em khác chỉ dùng nó để nhận đáp án. Vì vậy, công bằng trong thời đại AI không thể chỉ được đo bằng việc ``có quyền truy cập công cụ''; còn phải hỏi công cụ đang mở rộng hay thu hẹp cơ hội tham gia vào hoạt động trí tuệ có chất lượng.

Điều này đặc biệt đáng lưu ý khi vận dụng FHSM ở Việt Nam. Tiếp nhận FHSM nghiêm túc không có nghĩa đơn giản là đưa thêm bài khó. ``Khó'' và ``giàu suy lí'' là hai điều khác nhau. Một bài toán kĩ thuật rất khó vẫn có thể chỉ yêu cầu nhận dạng đúng dạng và thực hiện những thao tác đã luyện. Ngược lại, một bài toán dùng kiến thức không cao vẫn có thể buộc học sinh tìm cấu trúc, lựa chọn mô hình, biện minh và đánh giá nhiều cách tiếp cận. Vì vậy, công bằng về suy lí trước hết là công bằng về \emph{chất lượng hoạt động trí tuệ} mà mọi học sinh có cơ hội tham gia.

\phantomsection
\addcontentsline{toc}{section}{7. Từ chương trình chính thức tới lớp học: giá trị của FHSM như một khung trung gian}
\section*{7. Từ chương trình chính thức tới lớp học: giá trị của FHSM như một khung trung gian}

Đối sánh Chương trình giáo dục phổ thông môn Toán hiện hành với FHSM cho thấy hai văn bản tương thích đáng kể, nhưng phải hiểu đúng loại quan hệ giữa chúng. Chưa có cơ sở để nói Chương trình 2018 của Việt Nam được xây dựng từ FHSM hay chịu ảnh hưởng trực tiếp của cuốn sách. Hai văn bản cũng làm hai việc khác nhau. Chương trình quốc gia quy định mục tiêu, năng lực, nội dung và yêu cầu
cần đạt; FHSM tập trung vào chất lượng của hoạt động toán học mà học sinh tham gia. Vì vậy, điều đáng khảo sát là chúng \emph{hỗ trợ nhau được đến đâu}, chứ không phải văn bản nào là nguồn của văn bản nào.

Vùng giao nhau rõ nhất giữa hai văn bản nằm ở suy lí, nhưng không nên đồng nhất hai hệ khái niệm. Chương trình giáo dục phổ thông môn Toán 2018 xác định “tư duy và lập luận toán học” là một trong năm thành phần của năng lực toán học. Chương trình không chỉ nêu tên năng lực này mà còn mô tả những biểu hiện của nó: thực hiện các thao tác như so sánh, phân tích, tổng hợp, đặc biệt hóa, khái quát hóa, quy nạp và diễn dịch; sử dụng chứng cứ và lí lẽ; giải thích, chứng minh và điều chỉnh cách giải quyết vấn đề. Ở cấp trung học phổ thông, những yêu cầu ấy tiếp tục được nâng lên về mức độ thành thạo và độ phức tạp.

FHSM tiếp cận vấn đề từ một góc khác. Suy lí không được trình bày như một thành phần năng lực đứng bên cạnh các thành phần khác, mà như một quá trình nền tảng phải hiện diện trong việc học toán. Cuốn sách dành riêng các chương đầu để xác định suy lí và kiến tạo ý nghĩa, mô tả những thói quen suy lí, sự phát triển của chúng và cách thúc đẩy chúng trong lớp học; sau đó theo dõi hai quá trình này qua từng miền nội dung của toán trung học. Vì thế, giá trị bổ trợ của FHSM không nằm ở việc cung cấp những “biểu hiện” mà Chương trình 2018 hoàn toàn chưa có, mà ở việc triển khai dày hơn câu hỏi sư phạm: những hoạt động ấy xuất hiện ra sao trong một bài toán cụ thể, một miền nội dung cụ thể, một nhiệm vụ học tập cụ thể và trong tương tác giữa giáo viên với học sinh.

Cũng vì vậy, không nên hỏi đơn giản khái niệm nào “rộng hơn”. “Năng lực tư duy và lập luận toán học” của Chương trình 2018 là một năng lực cần hình thành ở người học; “suy lí” trong FHSM trước hết chỉ một quá trình nhận thức mà học sinh phải thường xuyên tham gia khi học toán. Hai cách tổ chức khái niệm khác nhau nhưng có vùng giao nhau lớn. FHSM có thể giúp giáo viên đọc những yêu cầu của chương trình dưới một câu hỏi cụ thể hơn: trong khi học một nội dung nhất định, học sinh đang được tạo cơ hội để quan sát, tìm quan hệ, hình thành và kiểm tra phỏng đoán, sử dụng căn cứ, giải thích, biện minh và xem lại cách nghĩ của mình như thế nào?

Kiến tạo ý nghĩa tạo ra một quan hệ bổ trợ khác. Chương trình Việt Nam không có một thành phần năng lực mang tên “kiến tạo ý nghĩa”, nhưng điều đó không có nghĩa quá trình này vắng mặt trong chương trình. Những hoạt động gần với nó xuất hiện ở nhiều nơi, chẳng hạn kết nối các ý tưởng toán học, mô hình hóa, giải quyết vấn đề, giao tiếp và vận dụng kiến thức trong những bối cảnh khác nhau. FHSM làm nổi bật một nguyên lí chung đứng sau những hoạt động ấy: người học phát triển sự hiểu biết về một tình huống, bối cảnh hay khái niệm bằng cách kết nối cái mới với tri thức đã có.

Bởi vậy, đóng góp có thể có của FHSM đối với Chương trình 2018 không phải là thay thế hay hoàn thiện một danh mục năng lực còn thiếu. Nó cung cấp một lăng kính sư phạm khác để xem xét cách những năng lực mà chương trình đã xác định được hình thành trong chính quá trình học nội dung toán học. Hai văn bản gặp nhau ở nhiều hoạt động, nhưng tổ chức chúng theo những cấu trúc khái niệm khác nhau; chính sự khác nhau ấy mới tạo ra khả năng bổ trợ.

Cách nhìn này đặc biệt hữu ích khi hai học sinh có thể đưa ra cùng một câu trả lời nhưng đã tham gia vào việc học theo những cách khác nhau. Chẳng hạn, cả hai đều có thể viết đúng một công thức; một em chỉ nhớ chuỗi kí hiệu, còn em kia có thể dựng lại công thức từ cấu trúc toán học và giải thích khi nào công thức được dùng. Nếu mục tiêu là phát triển năng lực, không thể chỉ dựa vào sự giống nhau của câu trả lời cuối cùng để coi hai trường hợp là tương đương. FHSM hướng sự chú ý trở lại quá trình học và những hoạt động toán học mà mỗi học sinh thực sự tham gia.

Sự tương thích còn xuất hiện ở giải quyết vấn đề và mô hình hóa. Chương trình Việt Nam xác định chúng như các thành phần năng lực riêng. FHSM đi sâu hơn vào những việc người học phải làm. Với giải quyết vấn đề, đó là phân tích tình huống, lựa chọn chiến lược, thực hiện có chủ đích, theo dõi tiến trình, diễn giải kết quả, kiểm tra tính hợp lí, xem lại giả định và khái quát hóa. Với mô hình hóa, FHSM nhấn mạnh việc đi qua lại giữa tình huống thực, mô hình toán học, các giả định, kết luận toán học, việc diễn giải và kiểm tra mô hình.

Cách mô tả này giúp tránh hai sự giản lược. ``Giải quyết vấn đề'' không chỉ là một loại bài tập mang nhãn riêng; ``mô hình hóa'' cũng không chỉ là viết một biểu thức từ dữ kiện.

Giao tiếp và công cụ cũng cho thấy quan hệ tương tự. Chương trình Việt Nam coi giao tiếp toán học và sử dụng công cụ, phương tiện học toán là thành phần năng lực. FHSM đặt giao tiếp trong hoạt động trình bày và thẩm định lí lẽ, còn công nghệ trong hoạt động tìm mẫu hình, tạo phỏng đoán, biểu diễn và giảm gánh nặng thao tác để dành nguồn lực cho suy nghĩ chiến lược. Trong kỉ nguyên AI, câu hỏi này trở nên sắc hơn: một công cụ đang mở rộng khả năng suy lí của người học hay đang thay thế chính hoạt động nhận thức mà bài học muốn hình thành?

Không có xung đột nền tảng giữa FHSM với việc Chương trình 2018 vẫn yêu cầu kiến thức và kĩ năng toán học cơ bản, thiết yếu. FHSM không kêu gọi bỏ công thức, thủ tục hay luyện tập. Nó yêu cầu học sinh không chỉ thực hiện chính xác mà còn hiểu vì sao thủ tục có hiệu lực, biết khi nào sử dụng và biết diễn giải kết quả. Xung đột chỉ xuất hiện nếu thủ tục trở thành mục đích tự thân, hoặc nếu ``khám phá'' được tổ chức mà không dẫn tới tri thức và năng lực toán học vững chắc.

Tuy nhiên, có tồn tại những căng thẳng thực hành. Một chương trình quốc gia phải bảo đảm phạm vi nội dung và yêu cầu cần đạt trong thời lượng hữu hạn; những nhiệm vụ giàu suy lí thường cần thời gian cho thử nghiệm, giải thích, so sánh, sửa sai và thảo luận. Nếu hiểu FHSM như một lớp hoạt động phải ``thêm vào'' sau khi đã dạy xong nội dung, nó dễ trở thành bất khả thi. Nhưng chính FHSM bác bỏ cách hiểu ấy: suy lí và kiến tạo ý nghĩa phải là cách học nội dung đang có, không phải một nội dung thứ sáu đứng cạnh đại số, hình học hay thống kê.

Căng thẳng thứ hai nằm ở đánh giá. Những điều chỉnh của Việt Nam từ năm 2025 đã hướng kì thi và đánh giá diện rộng về phía năng lực, nhưng một bài thi chuẩn hóa, thời gian hữu hạn không thể quan sát đầy đủ mọi tiến trình trí tuệ. Việc Bộ Giáo dục và Đào tạo tăng vai trò của đánh giá quá trình cũng phản ánh giới hạn này. FHSM vì thế không phủ định đánh giá chuẩn hóa; nó nhắc rằng những gì dễ đo trong một bài thi không nên bị đồng nhất với toàn bộ năng lực toán học mà chương trình muốn hình thành.

Từ những đối sánh trên, có thể tóm tắt quan hệ giữa hai văn bản bằng một công thức làm việc: \emph{Chương trình 2018 xác định đích đến; FHSM giúp làm rõ những hoạt động đưa người học tới đích đó}. Chương trình cho biết học sinh cần phát triển những năng lực nào và đạt những yêu cầu gì. FHSM giúp hỏi cụ thể hơn: nhiệm vụ nào, tương tác nào và cách biểu diễn nào đang buộc học sinh suy lí và kiến tạo ý nghĩa?

Trong môi trường AI, việc hiện thực hóa chương trình còn phải tính tới hai câu hỏi mới. Công cụ đang hỗ trợ hay làm thay những hoạt động nhận thức mà học sinh cần thực sự thực hiện? Và trong đánh giá, những biểu hiện quan sát được dưới điều kiện có công cụ cho phép ta suy ra điều gì về năng lực của chính người học? Theo nghĩa ấy, FHSM thích hợp nhất với Việt Nam khi được dùng như một \emph{khung trung gian để chất vấn và thiết kế cách chương trình được hiện thực hóa}, chứ không phải như một bộ chuẩn thay thế.

\phantomsection
\addcontentsline{toc}{section}{8. Ý nghĩa của bản chuyển ngữ: từ tiếp cận văn bản tới hình thành ngôn ngữ chuyên môn}
\section*{8. Ý nghĩa của bản chuyển ngữ: từ tiếp cận văn bản tới hình thành ngôn ngữ chuyên môn}

Giá trị đầu tiên của một bản dịch FHSM là giúp người đọc tiếp cận văn bản bằng tiếng Việt. Nhưng với một tác phẩm dựa nhiều vào những phân biệt khái niệm tinh tế, chuyển ngữ không thể chỉ là thay từ tiếng Anh bằng từ tiếng Việt. Nếu những khái niệm gần nhau trong nguyên tác bị nhập làm một, người đọc có thể mất chính những phân biệt mà lập luận của NCTM cần tới.

Quá trình hiệu chỉnh bản dịch này vì thế đã phải phân biệt có hệ thống \emph{reasoning}, \emph{inference}, \emph{argument}, \emph{argumentation}, \emph{justification}, \emph{proof} và \emph{explanation}; đồng thời tách \emph{sense making} khỏi \emph{understanding}, \emph{meaning} và \emph{interpretation}. Kết quả là một hệ tiếng Việt trong đó \emph{reasoning} được gọi là ``suy lí'', \emph{sense making} là ``kiến tạo ý nghĩa'', \emph{inference} là ``suy luận'', \emph{argument} là ``lập luận'', \emph{argumentation} là ``hoạt động lập luận'', \emph{justification} là ``biện minh'', \emph{proof} là ``chứng minh'', và \emph{interpretation} là ``diễn giải''. Lí do của từng lựa chọn được trình bày riêng trong phần \emph{Khảo luận về hệ thuật ngữ của bản chuyển ngữ} ở cuối sách; ở đây chỉ cần nhấn mạnh hệ quả học thuật của công việc ấy.

Ngôn ngữ chuyên môn không chỉ là những nhãn được gắn vào một tư tưởng đã có sẵn. Khi ngôn ngữ phân biệt được các khái niệm khác nhau, ta cũng có thể đặt những câu hỏi chính xác hơn. Chẳng hạn, nếu ``suy lí'', ``suy luận'' và ``lập luận'' bị nhập làm một, sẽ khó nói rõ học sinh đang làm gì: đang tham gia một hoạt động nhận thức rộng, đang suy ra một kết luận từ dữ liệu, hay đang trình bày một cấu trúc lí lẽ để người khác đánh giá. Tương tự, nếu ``kiến tạo ý nghĩa'' và ``sự hiểu biết'' đều bị rút về một chữ ``hiểu'', ta khó phân biệt quá trình tạo ra các kết nối với trạng thái hiểu biết được phát triển nhờ quá trình ấy.

Vì vậy, bản chuyển ngữ theo đuổi một mục tiêu khiêm tốn nhưng nền tảng: giúp các cuộc thảo luận bằng tiếng Việt giữ được những phân biệt khái niệm quan trọng của nguyên tác. Đây mới là mục tiêu của bản dịch, chưa phải một tác động xã hội đã được chứng minh. Mục tiêu ấy càng đáng chú ý khi Chương trình 2018 đã đưa ``tư duy và lập luận toán học'' vào cấu trúc năng lực, còn môi trường AI buộc giáo dục phải nói chính xác hơn về phần việc nhận thức nào người học cần tự thực hiện, khi nào công cụ được dùng để hỗ trợ, và dựa vào đâu ta có thể nhận định rằng hiểu biết hay năng lực đã thực sự hình thành ở người học.

Bản dịch đồng thời phải giữ ranh giới của mình. Nó không được dùng các quan niệm riêng của người chuyển ngữ để làm cho NCTM nói điều NCTM không nói. Chính vì vậy, phần bình giải được tách khỏi bản tiếng Việt và những khảo luận của người chuyển ngữ được ghi rõ tư cách văn bản. Nguyên tắc này đặc biệt quan trọng ở những vấn đề như trí tuệ nhân tạo hay sự mở rộng từ toán phổ thông lên toán đại học: chúng là những hướng đối thoại được gợi ra từ FHSM, không phải những phần đã có sẵn trong lập trường năm 2009 của NCTM.

\phantomsection
\addcontentsline{toc}{section}{9. Những giới hạn cần giữ khi tiếp nhận FHSM}
\section*{9. Những giới hạn cần giữ khi tiếp nhận FHSM}

Đánh giá cao một văn bản không có nghĩa xem nó như một kinh điển không thể chất vấn. Muốn sử dụng FHSM đúng, cần biết rõ cuốn sách dừng ở đâu.

Giới hạn thứ nhất là lịch sử và bối cảnh. FHSM được viết trong giáo dục Hoa Kì, gắn với những tranh luận về chương trình trung học, phân luồng, chuẩn, đánh giá, giáo dục đại học và nhu cầu lao động ở cuối thập niên 2000. Một số vấn đề có khả năng chuyển sang bối cảnh khác; một số khác gắn chặt với cơ cấu trường học Hoa Kì. Khi đọc ở Việt Nam, cần chuyển câu hỏi chứ không sao chép biểu hiện bề mặt. Chẳng hạn, phân luồng trong trường trung học Hoa Kì và những cơ chế phân hóa học sinh ở Việt Nam không thể mặc nhiên coi là cùng một hiện tượng, dù cả hai đều có thể được khảo sát bằng câu hỏi chung về cơ hội tiếp cận hoạt động toán học có chất lượng.

Giới hạn thứ hai là cấp học. FHSM tập trung vào trung học phổ thông. Nó có những phát biểu về sự liên thông từ các lớp trước tới giáo dục sau trung học, nhưng không cung cấp một lí thuyết đầy đủ về sự trưởng thành toán học ở đại học, về trừu tượng hóa cấu trúc, về chuyển tiếp sang chứng minh ở toán cao cấp, hay về hoạt động đặt vấn đề và kiến tạo lí thuyết trong nghiên cứu. Do đó, suy lí và kiến tạo ý nghĩa có thể là nền tảng quan trọng cho một tầm nhìn rộng hơn về giáo dục toán học, nhưng không nên tự động được tuyên bố là triết lí đầy đủ cho mọi trình độ.

Giới hạn thứ ba là bằng chứng. FHSM là một văn bản định hướng chính thức của NCTM, được xây dựng trên nghiên cứu và có Thư mục chú giải, nhưng bản thân nó không phải một tổng quan hệ thống theo chuẩn phương pháp luận hiện nay và không phải mọi khuyến nghị đều được kiểm chứng bằng cùng một loại bằng chứng. Tương tự, sự tương thích giữa FHSM với Chương trình 2018 mà tiểu luận này chỉ ra là một đối sánh khái niệm và cấu trúc; nó chưa phải bằng chứng thực nghiệm rằng triển khai FHSM ở Việt Nam sẽ tự động cải thiện kết quả học tập. Giá trị của sự đối sánh nằm ở việc tạo ra những giả thuyết thiết kế và câu hỏi nghiên cứu đáng kiểm định.

Giới hạn thứ tư liên quan trực tiếp đến trí tuệ nhân tạo. Mối liên hệ giữa FHSM với AI trong tiểu luận này là một \emph{diễn giải đương đại}. Nó được nâng đỡ bởi những vấn đề mà NCTM, UNESCO, các chính sách mới của Việt Nam và nghiên cứu gần đây đặt ra, nhưng không được dùng để viết lại lịch sử của FHSM. Hơn nữa, FHSM không đủ để trở thành một khung hoàn chỉnh cho giáo dục trong thời đại AI: các vấn đề về dữ liệu, quyền riêng tư, đạo đức, thiên lệch, quyền của người học, bất bình đẳng trong tiếp cận công cụ và quan hệ giữa năng lực chủ động của con người với máy cần những nguồn mới hơn. Giá trị của phép đọc nằm ở chỗ kiểm tra sức sống của hạt nhân suy lí và kiến tạo ý nghĩa dưới điều kiện mới, đồng thời nhận ra những tầng mà văn bản năm 2009 chưa thể giải quyết.

Những giới hạn ấy không làm FHSM mất giá trị. Ngược lại, chúng giúp xác định đúng cách sử dụng cuốn sách. FHSM có thể cung cấp những tiêu chuẩn và câu hỏi mạnh để suy nghĩ về giáo dục toán học, nhưng các tiêu chuẩn và câu hỏi đó luôn phải được đặt trong đối thoại với chương trình quốc gia, nghiên cứu mới, thực nghiệm tại địa phương và những biến đổi của xã hội.

\phantomsection
\addcontentsline{toc}{section}{10. Kết luận}
\section*{10. Kết luận}

Đóng góp trung tâm của \emph{Focus in High School Mathematics: Reasoning and Sense Making} có thể được nhìn như một sự chuyển hướng chú ý. Một hệ thống giáo dục thường dễ kiểm kê những thứ hữu hình: đã dạy chủ đề nào, dành bao nhiêu tiết, làm bao nhiêu bài tập, học sinh được bao nhiêu điểm và hoàn thành tới đâu. FHSM yêu cầu nhìn thêm vào một điều khó quan sát hơn: \emph{học sinh thực sự đang làm gì về mặt trí tuệ khi học toán}. Các em đang tìm quan hệ, phỏng đoán, biện minh, kiểm tra, giải thích và nối cái mới với cái đã biết như thế nào? Khi đặt câu hỏi ấy ở trung tâm, các vấn đề tưởng như tách rời---nhiệm vụ, chứng minh, công nghệ, đánh giá, công bằng và sự phối hợp của hệ thống---bắt đầu nối lại với nhau.

Đặt trong dòng phát triển của NCTM và bên cạnh những khung đương đại như PISA, FHSM cho thấy một sự nhấn mạnh mạnh mẽ vào suy lí và kiến tạo ý nghĩa ngay trong quá trình học nội dung toán học. Đặt trong bối cảnh Việt Nam, nó gặp một chương trình quốc gia đã xác định năm thành phần năng lực toán học và đang tiếp tục điều chỉnh đánh giá, năng lực số cùng giáo dục AI theo hướng phát triển năng lực. Vì vậy, giá trị của FHSM không nằm ở việc thuyết phục Việt Nam ``chuyển sang năng lực'', mà ở khả năng giúp đặt câu hỏi chi tiết hơn về cách những mục tiêu năng lực được hiện thực hóa thành nhiệm vụ, tương tác, biểu diễn, cách sử dụng công cụ và những hoạt động toán học mà học sinh thực sự tham gia trong mỗi miền nội dung.

Đặt trong kỉ nguyên trí tuệ nhân tạo, vấn đề tiến thêm một bước. AI có thể tham gia vào chính những phần việc mà một nhiệm vụ học tập vốn giao cho học sinh: đề xuất cách giải, xây dựng lập luận, viết lời giải thích, diễn giải hoặc kiểm tra kết quả. Vì thế, thiết kế dạy học phải xác định rõ công cụ đang hỗ trợ việc học ở đâu và ở đâu nó có nguy cơ làm thay hoạt động nhận thức mà người học cần thực sự trải qua.

Đánh giá cũng trở nên phức tạp hơn. Những gì học sinh nói, viết hay tạo ra trong một nhiệm vụ là bằng chứng để ta suy luận về hiểu biết và năng lực của các em, chứ không phải chính hiểu biết và năng lực ấy. Khi AI cùng tham gia tạo ra điều ta quan sát, điều kiện sử dụng công cụ phải trở thành một phần của việc diễn giải bằng chứng. Cùng một câu trả lời có thể mang ý nghĩa đánh giá khác nhau tùy vào việc học sinh tự thực hiện nhiệm vụ, được công cụ hỗ trợ một phần, hay đã giao cho công cụ phần việc mà ta đang muốn đánh giá.

Bởi thế, giá trị của FHSM đối với giáo dục toán học Việt Nam không nằm ở uy tín của một tổ chức Hoa Kì hay ở khả năng cung cấp một mô hình để sao chép. Giá trị ấy chỉ xuất hiện khi cuốn sách giúp chúng ta đặt được những câu hỏi tốt hơn đối với thực hành của chính mình. Một nhiệm vụ đang yêu cầu học sinh nghĩ gì? Các em có thực sự phải suy lí và kiến tạo ý nghĩa trên nội dung đang học hay không? Một công cụ đang mở rộng khả năng hoạt động toán học hay làm thay phần việc mà người học cần thực hiện? Một phép đánh giá đang thu thập bằng chứng gì, trong những điều kiện nào, và bằng chứng ấy cho phép kết luận tới đâu về năng lực? Học sinh có biết vì sao mình chấp nhận một kết quả do AI hỗ trợ tạo ra hay không? Và mọi nhóm học sinh có được tiếp cận những hoạt động toán học giàu suy lí như nhau hay không?

Nếu được đọc theo cách ấy, FHSM không cần trở thành một hệ triết học toàn năng để có giá trị lâu dài. Cuốn sách giúp đặt câu hỏi ``dạy toán gì'' bên cạnh câu hỏi ``học sinh đang làm gì về mặt trí tuệ khi học toán?'' Trong thời đại AI, ta còn phải hỏi thêm: công cụ đang tham gia vào quá trình học ở đâu, phần hoạt động nào vẫn cần thuộc về người học, và dựa vào bằng chứng nào ta có thể nói rằng năng lực và sự hiểu biết đã thực sự hình thành ở chính các em? Những câu hỏi này không thay thế nhau. Chúng bổ sung cho nhau để xác định khi nào một nội dung toán học thực sự trở thành năng lực và sự hiểu biết của người học.

\phantomsection
\addcontentsline{toc}{section}{Tài liệu tham khảo của tiểu luận}
\section*{Tài liệu tham khảo của tiểu luận}

\begingroup
\small
\setlength{\parindent}{0pt}
\setlength{\parskip}{0.55\baselineskip}

Bastani, H., Bastani, O., Sungu, A., Ge, H., Kabakcı, Ö., và Mariman, R.
(2025).
``Generative AI without guardrails can harm learning: Evidence from high
school mathematics.''
\emph{Proceedings of the National Academy of Sciences}, 122(26),
e2422633122.
doi: 10.1073/pnas.2422633122.

Bộ Giáo dục và Đào tạo (2018).
\emph{Chương trình giáo dục phổ thông}, ban hành kèm theo Thông tư
32/2018/TT-BGDĐT ngày 26 tháng 12 năm 2018, cùng các văn bản sửa đổi, hợp nhất
hiện hành.

Bộ Giáo dục và Đào tạo (2025).
\emph{Khung năng lực số cho người học}, ban hành kèm theo Thông tư
02/2025/TT-BGDĐT ngày 24 tháng 1 năm 2025.

Bộ Giáo dục và Đào tạo (2026a).
\emph{Quy định về đánh giá diện rộng cấp quốc gia chất lượng giáo dục phổ
thông}, ban hành kèm theo Thông tư 08/2026/TT-BGDĐT ngày 26 tháng 2 năm 2026.

Bộ Giáo dục và Đào tạo (2026b).
\emph{Khung nội dung giáo dục trí tuệ nhân tạo (AI) cho học sinh phổ thông}, ban
hành kèm theo Quyết định 2422/QĐ-BGDĐT ngày 18 tháng 8 năm 2026.

Bộ Giáo dục và Đào tạo (2026c).
``Hướng dẫn triển khai thực hiện nội dung giáo dục trí tuệ nhân tạo cho học
sinh phổ thông từ năm học 2026--2027'', Công văn 5588/BGDĐT-GDPT ngày 19
tháng 8 năm 2026.

Cục Quản lí chất lượng, Bộ Giáo dục và Đào tạo (2023).
``Cấu trúc định dạng đề thi tốt nghiệp THPT từ năm 2025'', công bố ngày
29 tháng 12 năm 2023.

Kilpatrick, J., Swafford, J., và Findell, B. (chủ biên) (2001).
\emph{Adding It Up: Helping Children Learn Mathematics}.
Washington, DC: National Academy Press.

National Research Council (2001).
\emph{Knowing What Students Know: The Science and Design of Educational Assessment}.
Washington, DC: National Academy Press.
doi: 10.17226/10019.

Miao, F., và Holmes, W. (2023).
\emph{Guidance for Generative AI in Education and Research}.
Paris: UNESCO.

Miao, F., Shiohira, K., và Lao, N. (2024).
\emph{AI Competency Framework for Students}.
Paris: UNESCO.
doi: 10.54675/JKJB9835.

National Council of Teachers of Mathematics (NCTM) (1980).
\emph{An Agenda for Action: Recommendations for School Mathematics of the
1980s}. Reston, VA: NCTM.

National Council of Teachers of Mathematics (NCTM) (1989).
\emph{Curriculum and Evaluation Standards for School Mathematics}.
Reston, VA: NCTM.

National Council of Teachers of Mathematics (NCTM) (2000).
\emph{Principles and Standards for School Mathematics}.
Reston, VA: NCTM.

National Council of Teachers of Mathematics (NCTM) (2006).
\emph{Curriculum Focal Points for Prekindergarten through Grade 8
Mathematics: A Quest for Coherence}.
Reston, VA: NCTM.

National Council of Teachers of Mathematics (NCTM) (2009).
\emph{Focus in High School Mathematics: Reasoning and Sense Making}.
Reston, VA: NCTM.

National Council of Teachers of Mathematics (NCTM) (2024).
\emph{Artificial Intelligence and Mathematics Teaching}.
Position Statement, February 2024. Reston, VA: NCTM.

OECD (2023).
\emph{PISA 2022 Assessment and Analytical Framework}.
Paris: OECD Publishing.
doi: 10.1787/dfe0bf9c-en.

UNESCO (2025a).
\emph{AI and the Future of Education: Disruptions, Dilemmas and Directions}.
Paris: UNESCO.

UNESCO (2025b).
\emph{AI and Education: Protecting the Rights of Learners}.
Paris: UNESCO.

\endgroup

\endgroup
